\subsection{3D Gaussian Splatting Primer}

3D Gaussian Splatting~\cite{kerbl2023_3dgs} represents a 3D scene as a set
of $N$ Gaussian primitives, each parameterized by:
\begin{itemize}
    \item Position $\boldsymbol{\mu}_i \in \mathbb{R}^3$ (center)
    \item Covariance $\boldsymbol{\Sigma}_i = \boldsymbol{R}_i
    \boldsymbol{S}_i \boldsymbol{S}_i^T \boldsymbol{R}_i^T$ (decomposed into
    rotation $\boldsymbol{R}_i$ and scale $\boldsymbol{S}_i$ for positive
    semi-definiteness)
    \item Opacity $\alpha_i \in [0,1]$
    \item Color via spherical harmonics (SH) coefficients up to degree~3
\end{itemize}
Rendering is performed via differentiable tile-based rasterization: each
pixel's color is an alpha-composited blend of projected Gaussians,
\begin{equation}
    \boldsymbol{C}(\boldsymbol{p}) = \sum_{i \in \mathcal{N}}
    \boldsymbol{c}_i \, \alpha_i \prod_{j=1}^{i-1}(1-\alpha_j),
\end{equation}
where $\alpha_i$ is the product of the 2D-projected Gaussian opacity and
the per-pixel opacity.  Depth rendering is similarly obtained by
alpha-compositing the Gaussian centers along each ray.

\subsection{From NeRF-SLAM to Gaussian-SLAM}

Early radiance-field SLAM systems used NeRF-based implicit
representations~\cite{sucar2021_imap,zhu2022_niceslam,wang2023_coslam,
rosinol2023_nerfslam,sandstrom2023_pointslam}.  While achieving high-quality
rendering, these methods were limited by slow volumetric rendering and
MLP-capacity bottlenecks.  The transition to 3DGS brought three key
advantages: (1)~real-time rasterization ($>$100~FPS), (2)~explicit
primitives enabling direct geometric operations (ICP, nearest-neighbor
queries), and (3)~incremental map expansion without retraining.

\subsection{The Depth-Source Challenge}

In a typical 3DGS-SLAM pipeline, the depth source serves three roles:
\begin{enumerate}
    \item \textbf{Initialization}: Each Gaussian's center $\boldsymbol{\mu}$
    is initialized by back-projecting depth pixels into 3D world coordinates.
    Sparse or inaccurate depth leads to under-reconstruction.
    \item \textbf{Supervision}: Depth rendering loss
    $\mathcal{L}_{\text{depth}} = \|\hat{D} - D\|_1$ constrains the Gaussian
    geometry.  Without reliable depth, Gaussians may converge to incorrect
    positions that overfit training views.
    \item \textbf{Tracking}: In map-centric SLAM (e.g.,
    MonoGS~\cite{matsuki2024monogs}), rendered depth is compared against
    sensor depth to estimate camera pose.  Depth quality directly impacts
    tracking accuracy.
\end{enumerate}

This depth-source dependency is the central axis along which we organize
this survey.  Table~\ref{tab:depth_sources} summarizes the three modalities.

\begin{table}[t]
\centering
\caption{Comparison of depth source modalities for Gaussian Splatting SLAM.}
\label{tab:depth_sources}
\small
\begin{tabular}{lcccc}
\toprule
\textbf{Modality} & \textbf{Metric Scale} & \textbf{Range} & \textbf{Cost} & \textbf{Sparsity} \\
\midrule
LiDAR + RGB + IMU & Yes & Long ($>$100\,m) & High & Sparse \\
iToF + RGB + IMU & Yes & Short ($<$10\,m) & Low & Very Sparse \\
Mono Depth + RGB + IMU & Via IMU & Unlimited & Minimal & Dense (learned) \\
\bottomrule
\end{tabular}
\end{table}

\subsection{Evaluation Metrics}

Gaussian SLAM systems are typically evaluated on:
\begin{itemize}
    \item \textbf{Tracking}: Absolute Trajectory Error (ATE) in cm or m.
    \item \textbf{Rendering}: PSNR (dB), SSIM, LPIPS on held-out views.
    \item \textbf{Geometry}: Depth L1 error, mesh accuracy/completeness
    (cm), Chamfer distance.
    \item \textbf{Efficiency}: FPS, GPU memory (GB), CPU utilization.
\end{itemize}
Common datasets include Replica (indoor), TUM RGB-D, ScanNet, ETH3D, KITTI
(outdoor driving), and the MARS-LVIG dataset for LiDAR-visual-inertial
sequences.
